\chapter{Results}
\label{ch:results}

\section{The evaluation protocol}

Four variants, sharing every stage before the branch, are run end to end on
seven real plans that no model has seen, under both sets of score weights.
The plans are drawn from Swiss Dwellings, the only source where the
labelling procedure applies exactly as written, and were selected to be
typical rather than favourable: each plan in the corpus was scored by its
distance from the corpus median across seven scale-independent
measurements, the 120 closest were sorted by size, and seven were taken at
even intervals so that the set spans small to large.
Table~\ref{tab:demo} gives them.

\begin{table}[H]
\centering
\caption[The seven evaluation plans]{The seven evaluation plans. Marked
columns are those the original drawing recorded, which the generator
retains in every candidate. The final column is the rank the derived
arrangement receives from the baseline variant under learned weights.}
\label{tab:demo}
\small
\begin{tabular}{llrrrrrr}
\toprule
Plan & Source id & Footprint (m) & Grid & Cols. & Marked & Candidates & Rank (V3) \\
\midrule
demo\_01 & 5855 & 12.5$\times$9.4 & 9$\times$8 & 43 & 1 & 3,963 & 5 \\
demo\_02 & 8491 & 15.5$\times$13.6 & 10$\times$9 & 72 & 0 & 3,999 & 9 \\
demo\_03 & 44116 & 24.3$\times$15.4 & 10$\times$9 & 55 & 0 & 3,996 & 61 \\
demo\_04 & 16422 & 19.8$\times$16.5 & 9$\times$8 & 42 & 8 & 512 & 1 \\
demo\_05 & 2001 & 17.8$\times$19.1 & 9$\times$10 & 65 & 0 & 4,000 & 41 \\
demo\_06 & 41319 & 24.4$\times$17.3 & 10$\times$8 & 68 & 0 & 3,935 & 2 \\
demo\_07 & 840 & 36.3$\times$37.7 & 10$\times$11 & 56 & 0 & 3,975 & 33 \\
\bottomrule
\end{tabular}
\end{table}

For each plan the pipeline parses the walls, enumerates candidates,
re-derives the label from that plan's own walls, and runs each variant. The
measurements taken are whether the correct arrangement is the variant's
first choice and whether it is among the first five; the rank of the
correct arrangement in the variant's full ordering, which is more
informative than a hit rate when hit rates are low; and two measures of
near-agreement, namely the average distance from each chosen column to the
nearest correct one and the proportion of columns that coincide within a
tolerance. Requiring an exact match alone would discard the distinction
between a near miss and a complete failure.

Variants are compared by paired tests across the seven plans, a Wilcoxon
signed-rank test together with an exact sign test, rather than by placing
two averages side by side. Seven plans provide limited statistical power,
which is stated with the results rather than left for the reader to infer.

\section{Candidate coverage}
\label{sec:coverage}

Coverage is reported before any accuracy figure because it caps all of
them. If the correct arrangement is absent from the candidate set, an
accuracy figure measures that absence rather than the quality of any
variant, and all four variants would be measuring the same ceiling.

\textbf{The derived arrangement is present in the candidate set for all
seven plans.} Reaching that required five changes, each documented in the
stage it affects: aligning the parser with the labelling procedure as
Section~\ref{sec:parser} describes; adding the wall-based placement rule to
the generator; treating marked columns as fixed input; widening the
permitted span range so that the correct arrangement's own grid lines are
not excluded in advance, since load-bearing walls in these dwellings can
stand 0.13\,m apart; and excluding angled-wall positions, which the
labelling procedure cannot represent. The accuracy figures that follow are
therefore unconditional.

Candidate set sizes range from 512 to 4,000 per plan, the upper figure
being the configured limit at which the generator begins sampling.

\section{The judges on their own task}
\label{sec:mlresults}

Before the variants are compared end to end, each judge is measured on the
classification task it was trained for.

\subsection{The feature judge}

Table~\ref{tab:ml1} reports the feature judge. The headline figure is
0.862 ROC AUC over five grouped splits.

\begin{table}[H]
\centering
\caption[The feature judge]{The feature judge. The pooled figure is the
mean over five grouped splits with its standard deviation; the remaining
rows are single deterministic runs. ``Against transplant examples only''
restricts the negative class to one kind at a time.}
\label{tab:ml1}
\small
\begin{tabular}{lrr}
\toprule
Experiment & ROC AUC & PR AUC \\
\midrule
Feature judge, five grouped splits & 0.862 $\pm$ 0.006 & 0.831 $\pm$ 0.008 \\
Best single-measurement threshold & 0.679 $\pm$ 0.008 & --- \\
\midrule
\quad tested on Swiss only & 0.896 & 0.865 \\
\quad tested on CubiCasa only & 0.774 & 0.747 \\
\midrule
Trained Swiss, tested CubiCasa & 0.673 & 0.608 \\
Trained CubiCasa, tested Swiss & 0.805 & 0.727 \\
\midrule
Against transplant examples only & 0.668 & 0.884 \\
Against jitter examples only & 0.960 & 0.989 \\
Against shear examples only & 0.956 & 0.984 \\
Against dropout examples only & 0.881 & 0.965 \\
\bottomrule
\end{tabular}
\end{table}

Four observations follow.

\textbf{The reported figure fell from 0.947 to 0.862 when the training
examples were corrected}, as Section~\ref{sec:negatives} described.
Approximately a tenth of the previous figure was the model recognising its
own example generator. The reduction is reported as a finding rather than
adjusted away, and it is the clearest single illustration of why the
construction of negative examples deserves attention.

\textbf{No single measurement approaches the full model.} The best
single-threshold rule reaches 0.679, so the remaining signal genuinely
requires several measurements in combination.

\textbf{Transplanted arrangements are the hardest class} at 0.668, since
they are real arrangements merely placed on the wrong building. They are
also the class closest to the judge's actual task in the pipeline, where
every candidate is geometrically valid and the question is which suits
this plan, so the pooled figure should be read alongside them.

\textbf{Transfer between sources is asymmetric.} A model trained on Swiss
drawings and tested on CubiCasa reaches 0.673, while the reverse reaches
0.805. Cross-source claims therefore rest on the pooled model, and the
asymmetry is reported rather than averaged out.

The two classes are balanced in equal proportion by construction, so
neither figure should be read as an estimate of performance at deployment
prevalence; both ROC AUC and PR AUC are reported for completeness.

\subsection{The image judge, and what its advantage consists of}

Table~\ref{tab:ml2} reports the image judge together with the two controls
that interpret it: the same network with its wall layer blanked, and the
feature judge fitted on the measurements of exactly the same rows and
scored on exactly the same splits.

\begin{table}[H]
\centering
\caption[The image judge and its controls]{The image judge and its
controls, over three grouped splits. The feature judge here is
\emph{paired}: fitted on the measurements of the same rows and scored on
the same split indices, which removes the confound present in an earlier
version of this comparison.}
\label{tab:ml2}
\small
\begin{tabular}{lrr}
\toprule
Model & ROC AUC & PR AUC \\
\midrule
Image judge, walls and columns & 0.995 $\pm$ 0.003 & 0.995 $\pm$ 0.003 \\
Image judge, walls layer blanked & 0.874 $\pm$ 0.019 & 0.819 $\pm$ 0.032 \\
Feature judge, paired on the same rows & 0.860 $\pm$ 0.007 & 0.830 $\pm$ 0.010 \\
\midrule
Image judge, walls and columns, Swiss$\rightarrow$CubiCasa & 0.983 & 0.983 \\
Feature judge, paired on the same rows, Swiss$\rightarrow$CubiCasa & 0.673 & 0.608 \\
Image judge, walls and columns, CubiCasa$\rightarrow$Swiss & 0.988 & 0.989 \\
Feature judge, paired on the same rows, CubiCasa$\rightarrow$Swiss & 0.805 & 0.727 \\
Image judge, walls layer blanked, Swiss$\rightarrow$CubiCasa & 0.771 & 0.654 \\
Image judge, walls layer blanked, CubiCasa$\rightarrow$Swiss & 0.769 & 0.699 \\
\bottomrule
\end{tabular}
\end{table}

Figure~\ref{fig:cnnvsfeatures} shows the three side by side, and the middle
bar carries the finding.

\begin{figure}[H]
\centering
\includegraphics[width=0.75\textwidth]{11_cnn_vs_features}
\caption[Where the image judge's advantage comes from]{Where the image
judge's advantage comes from. Given the same information as the feature
judge, that is columns without walls, the network performs equivalently.
The advantage appears only when the wall layer is present.}
\label{fig:cnnvsfeatures}
\end{figure}

\textbf{Given the same information, the two judges are
indistinguishable}: 0.874 for the picture against 0.860 for the
measurements, with overlapping intervals. \textbf{Given the walls, the
picture is far better} at 0.995, and the advantage survives testing on a
source the model never saw, at 0.983 and 0.988.

The conclusion is therefore about information rather than architecture. The
convolutional network does not read column patterns better than a
gradient-boosted model reads measurements of them. The picture wins because
it can show whether the columns coincide with the walls, and the
scale-independent measurements cannot carry that relationship across
sources drawn in different units.

One qualification travels with this result wherever it appears. The derived
labels place columns inside load-bearing walls \emph{by definition}, so the
strength of the column-wall registration signal is partly a restatement of
the labelling procedure. The same qualification applies to the off-wall
measure and to its large learned weight in Table~\ref{tab:weights}. These
are treated as one finding with three appearances: on this corpus, the
strongest available indication that an arrangement is real is whether its
columns coincide with its walls, and every representation permitted to see
the walls discovers it.

\subsection{The quantum judge}

Table~\ref{tab:quantum} reports the size sweep, and
Figure~\ref{fig:quantum} plots it.

\begin{table}[H]
\centering
\caption[The quantum judge across training-set sizes]{The quantum judge
across training-set sizes, all models scored on one fixed held-out test
set. Q marks quantum models and C classical. ``Gain'' is the change from
the smallest to the largest training set; ``Time'' is the fitting time at
the largest.}
\label{tab:quantum}
\scriptsize
\begin{tabular}{llrrrrrr}
\toprule
Model & & \multicolumn{4}{c}{Training rows} & Gain & Time (s) \\
\cmidrule(lr){3-6}
 & & 400 & 2,000 & 8,000 & 19,936 & & \\
\midrule
Boosting, all 6 measurements & C & 0.790 & 0.841 & 0.860 & 0.865 & $+0.075$ & 5 \\
Boosting, same 4 measurements & C & 0.705 & 0.753 & 0.779 & 0.789 & $+0.084$ & 3 \\
Logistic, same 4 (capacity-matched) & C & 0.747 & 0.749 & 0.746 & 0.746 & $-0.001$ & 0 \\
\midrule
VQC, one encoding, 2 layers & Q & 0.741 & 0.751 & 0.734 & 0.748 & $+0.007$ & 40 \\
Data re-uploading, 6 layers & Q & 0.741 & 0.761 & 0.756 & 0.746 & $+0.005$ & 251 \\
\bottomrule
\end{tabular}
\end{table}

\begin{figure}[H]
\centering
\includegraphics[width=0.85\textwidth]{10_quantum_vs_classical}
\caption[Quantum and classical models against training-set size]{Quantum
and classical models against training-set size, on one fixed test set. The
circuits track the capacity-matched linear control; only the tree
ensembles convert additional data into accuracy.}
\label{fig:quantum}
\end{figure}

Three findings follow, and the third is the one worth stating in a
defence.

\textbf{Given the same information, the circuits match their
capacity-matched control.} At every size the two quantum models lie within
about 0.015 of logistic regression on the same four measurements, and at
2,000 rows they are slightly ahead. A four-qubit circuit can learn this
task.

\textbf{The circuits gain almost nothing from more data.} Fifty times more
training data moves them by 0.007 and 0.005, which is within run-to-run
variation. The capacity-matched linear control behaves identically at
$-0.001$, while the tree ensembles convert the same additional data into
0.075 and 0.084.

\textbf{The ceiling is therefore capacity, not data.} Because the circuits
improve with data in the manner of a linear model rather than a flexible
one, the limit on this task lies in the expressive capacity of four qubits
and two entangling layers, not in the number of examples available. That
distinction cannot be drawn from a single small-sample run, which is why
the experiment sweeps the training size instead of reporting one figure.

The difference in computational cost is part of the result. At the largest
size the plain circuit required 40 seconds of simulation and the
re-uploading circuit 251, against 3 to 5 seconds for the tree ensembles and
under a second for logistic regression, to reach a lower score.

\subsection{A cautionary result retained deliberately}

On an earlier version of the labels, which placed a column at every grid
crossing, this stage appeared to show a quantum ensemble matching a strong
classical model. When the labels were corrected the finding reversed
completely and the ensemble became the weakest model in the table. The
circuits had not changed; the task had. The episode is retained because it
demonstrates how readily comparisons of this kind produce misleadingly
favourable results, and it is the reason the controls, rather than the
circuits, constitute this stage's contribution.

\section{The score weights}
\label{sec:weightresults}

Table~\ref{tab:weights} compares the two weight vectors.

\begin{table}[H]
\centering
\caption[Fixed and learned score weights]{Fixed and learned score weights.
The learned vector is normalised to a mean absolute weight of one, the same
scale as the fixed vector; ranking is invariant to that scale.}
\label{tab:weights}
\small
\begin{tabular}{lrrr}
\toprule
Measure & Fixed weight & Learned weight & S.d.\ over splits \\
\midrule
Density & 1.00 & $-0.42$ & 0.03 \\
Span demand & 1.00 & $+0.38$ & 0.17 \\
Irregularity & 1.00 & $-0.46$ & 0.06 \\
Off-wall & 1.00 & $+2.77$ & 0.09 \\
Repetition & 1.00 & $+0.97$ & 0.01 \\
\bottomrule
\end{tabular}
\end{table}

On held-out plans the learned weights rank the real arrangement above an
alternative in 89.4 per cent of paired comparisons, with a standard
deviation of 0.6 percentage points across splits, against 74.7 per cent for
the fixed weights on the same data.

The signs are the informative part. Span demand and repetition carry the
signs that design reasoning predicts: real arrangements have shorter loaded
spans and more repeated bays than the alternatives generated on their own
grids. The off-wall measure carries the largest weight, though partly by
construction as noted above. Density and irregularity, however, carry
\emph{negative} weights, which means the corpus contains real arrangements
that are \emph{denser} and \emph{less regular} than the alternatives.

That is not an error but a property of the buildings. These are
wall-bearing dwellings in which columns stand where the load-bearing walls
stand, and those walls follow rooms rather than a regular structural
module. Figure~\ref{fig:weights} shows the two vectors, and this single
result predicts the central pattern of the accuracy table that follows.

\begin{figure}[H]
\centering
\includegraphics[width=0.78\textwidth]{06_weights}
\caption[The two sets of score weights]{The two sets of score weights. The
learned vector inverts the sign on density and irregularity relative to
the design intuition that the fixed vector encodes, because real
wall-derived arrangements are denser and less regular than the
alternatives available on the same grid.}
\label{fig:weights}
\end{figure}

\section{The four variants end to end}
\label{sec:variants}

Table~\ref{tab:accuracy} is the principal result of this thesis.

\begin{table}[H]
\centering
\caption[Accuracy of the four variants]{Accuracy of the four variants under
both sets of score weights, averaged over the seven evaluation plans.
Rank is the position of the derived arrangement in the variant's full
ordering of 512 to 4,000 candidates. Distance is the mean Chamfer distance
between chosen and correct columns; overlap is the proportion coinciding
within tolerance.}
\label{tab:accuracy}
\small
\begin{tabular}{lrrrrrr}
\toprule
Variant & Mean rank & Median rank & Top-1 & Top-5 & Distance (m) & Overlap \\
\midrule
\multicolumn{7}{l}{\emph{Fixed weights}} \\
\quad V1 (feature judge) & 2429.3 & 2395 & 0.000 & 0.000 & 1.000 & 0.287 \\
\quad V2 (image judge) & 2429.3 & 2395 & 0.000 & 0.000 & 0.958 & 0.312 \\
\quad V3 (score only, baseline) & 2429.3 & 2395 & 0.000 & 0.000 & 1.000 & 0.287 \\
\quad V4 (quantum judge) & 2429.3 & 2395 & 0.000 & 0.000 & 1.000 & 0.287 \\
\midrule
\multicolumn{7}{l}{\emph{Learned weights}} \\
\quad V1 (feature judge) & 21.9 & 10 & 0.000 & 0.429 & 0.248 & 0.704 \\
\quad V2 (image judge) & 21.7 & 9 & 0.143 & 0.429 & 0.264 & 0.718 \\
\quad V3 (score only, baseline) & 21.7 & 9 & 0.143 & 0.429 & 0.271 & 0.722 \\
\quad V4 (quantum judge) & 22.0 & 11 & 0.143 & 0.429 & 0.282 & 0.683 \\
\bottomrule
\end{tabular}
\end{table}

\subsection{The score weights dominate}

Under fixed weights the derived arrangement sits at a median rank of about
2,395 out of roughly 4,000 for every variant, and no variant recovers it
first on any plan. The equal-weight score prefers sparse, regular grids,
while real wall-derived arrangements are dense and irregular; the negative
density and irregularity coefficients of
Section~\ref{sec:weightresults} predicted exactly this. With the correct
arrangement never reaching the shortlist the judges see, a judge can only
change which \emph{wrong} answer is selected.

Under learned weights the picture changes entirely. The score \emph{alone}
places the correct arrangement at a median rank of 9, recovers it exactly
on one of the seven plans, and averages 0.271\,m Chamfer distance with 0.722
overlap. The difference between the two halves of
Table~\ref{tab:accuracy} is far larger than any difference within either
half: \textbf{the choice of weights matters much more than the choice of
judge}.

\subsection{The judges are second-order}

Within the learned half, the image judge trims the Chamfer distance
slightly, from 0.271 to 0.264\,m, but also lowers the overlap slightly,
from 0.722 to 0.718: a mixed result, not a consistent improvement. The
feature judge trades a first-place recovery away for the best distance of
the four (0.248\,m), and the quantum judge, now trained on the same 80\%
training partition as the other two, recovers the arrangement first on
the same one plan as the baseline and the image judge, and returns the
weakest agreement of the four. Figure~\ref{fig:accuracy} shows both
halves.

\begin{figure}[H]
\centering
\includegraphics[width=\textwidth]{12_accuracy}
\caption[Geometric agreement by variant and weight set]{Geometric
agreement by variant and weight set. Pale bars are fixed weights and solid
bars learned weights. The gap between the two weight sets exceeds every
difference between variants within either set.}
\label{fig:accuracy}
\end{figure}

\subsection{Significance}

Table~\ref{tab:tests} gives the paired tests for the learned half, which is
the only half in which the variants differ at all.

\begin{table}[H]
\centering
\caption[Paired significance tests]{Paired significance tests across the
seven evaluation plans, learned weights, for every comparison in which the
two variants differ. No comparison reaches conventional significance.}
\label{tab:tests}
\small
\begin{tabular}{llrrrr}
\toprule
Comparison & Metric & Mean A & Mean B & Wilcoxon $p$ & Sign $p$ \\
\midrule

\bottomrule
\end{tabular}
\end{table}

\textbf{No variant separates from the baseline at conventional
significance.} Every paired Wilcoxon test returns $p = 1.000$: each pair
of variants differs on at most two of the seven plans, and the
differences run in both directions. Any difference between variants is
well within the noise at this sample size.

This is the answer to the third research question of
Section~\ref{sec:questions}, and it is reported as the outcome of the
experiment rather than presented as a shortfall. The protocol was designed
so that this result would be publishable: the baseline was retained, the
comparison was paired, and the sample size and its consequences are
stated.

\section{Sensitivity analyses}
\label{sec:sensitivity}

Every hand-chosen constant in the pipeline is varied and its influence
measured.

\subsection{The score weights}

Table~\ref{tab:sensitivity} perturbs each weight by $\pm 50$ per cent and
removes each measure in turn.

\begin{table}[H]
\centering
\caption[Sensitivity to the score weights]{Sensitivity of the baseline
variant to the score weights, fixed-weight mode. Kendall's $\tau$ compares
the perturbed ranking with the nominal one; top-5 overlap is the
proportion of the nominal top five that survives.}
\label{tab:sensitivity}
\small
\setlength{\tabcolsep}{5pt}
\begin{tabular}{@{}lrrrr@{}}
\toprule
Perturbation & Kendall $\tau$ & Top-1 kept & Top-5 kept & Rank \\
\midrule
Nominal (unperturbed) & 1.000 & 1.000 & 1.000 & 2429 \\
Density, +50% & 0.881 & 0.857 & 0.857 & 2873 \\
Density, -50% & 0.853 & 0.857 & 0.943 & 1556 \\
Density, removed & 0.690 & 0.571 & 0.829 & 643 \\
Irregularity, +50% & 0.902 & 0.714 & 0.914 & 2546 \\
Irregularity, -50% & 0.870 & 0.714 & 0.829 & 2236 \\
Irregularity, removed & 0.699 & 0.143 & 0.543 & 2002 \\
Off-wall, +50% & 0.915 & 1.000 & 1.000 & 2213 \\
Off-wall, -50% & 0.896 & 0.857 & 0.886 & 2660 \\
Off-wall, removed & 0.778 & 0.857 & 0.771 & 2861 \\
Repetition, +50% & 0.901 & 0.857 & 0.886 & 2138 \\
Repetition, -50% & 0.875 & 0.571 & 0.800 & 2658 \\
Repetition, removed & 0.725 & 0.429 & 0.486 & 2838 \\
Span demand, +50% & 0.865 & 0.571 & 0.857 & 2122 \\
Span demand, -50% & 0.851 & 1.000 & 0.829 & 2702 \\
Span demand, removed & 0.704 & 0.429 & 0.629 & 2888 \\
\bottomrule
\end{tabular}
\end{table}

Varying any single weight by half keeps Kendall's $\tau$ between 0.85 and
0.92 and leaves four-fifths of the top five unchanged. That is the
empirical defence of equal weighting promised in
Section~\ref{sec:weights}: the vector cannot be argued for, but the
ranking barely depends on it. Removing a measure entirely costs more, with
$\tau$ falling to between 0.69 and 0.78, so every measure carries real
ordering information and none is decoration.

One perturbation is diagnostic. Removing the density measure
\emph{improves} the correct arrangement's mean rank from 2,429 to 643,
which independently confirms from a different direction the negative
density coefficient that the learned weights found.

\subsection{The judge's authority}

Table~\ref{tab:lambda} varies $\lambda$ from 0 to 1.

\begin{table}[H]
\centering
\caption[Sensitivity to the judge's weight]{Sensitivity to $\lambda$, the
weight given to the judge in Equation~\eqref{eq:blend}, for variant V1. At
$\lambda = 0$ every variant reduces to the baseline.}
\label{tab:lambda}
\small
\begin{tabular}{lrrr}
\toprule
Judge weight & Mean rank & Top-1 & Top-5 \\
\midrule
\multicolumn{4}{l}{\emph{Fixed weights}} \\
\quad $\lambda = 0.0$ & 2429.3 & 0.000 & 0.000 \\
\quad $\lambda = 0.1$ & 2429.3 & 0.000 & 0.000 \\
\quad $\lambda = 0.3$ & 2429.3 & 0.000 & 0.000 \\
\quad $\lambda = 0.5$ & 2429.3 & 0.000 & 0.000 \\
\quad $\lambda = 1.0$ & 2429.3 & 0.000 & 0.000 \\
\midrule
\multicolumn{4}{l}{\emph{Learned weights}} \\
\quad $\lambda = 0.0$ & 21.7 & 0.143 & 0.429 \\
\quad $\lambda = 0.1$ & 21.6 & 0.143 & 0.429 \\
\quad $\lambda = 0.3$ & 21.9 & 0.000 & 0.429 \\
\quad $\lambda = 0.5$ & 21.9 & 0.000 & 0.429 \\
\quad $\lambda = 1.0$ & 21.9 & 0.000 & 0.429 \\
\bottomrule
\end{tabular}
\end{table}

In fixed-weight mode $\lambda$ is inert, for the reason given above: the
correct arrangement never reaches the shortlist the judge sees. In learned
mode the effect is real, small and mixed in both directions. Raising
$\lambda$ from 0 to 1 worsens the mean rank slightly, from 21.7 to 21.9,
costs one first-place recovery, and gains no top-five placing. The default
of 0.3 is defensible precisely because so little depends on it.

\subsection{The labelling threshold}

Table~\ref{tab:thresholdacc} is the most consequential sensitivity analysis
in this report, because every figure in it rests on the 160\,mm constant.

\begin{table}[H]
\centering
\caption[Sensitivity to the labelling threshold]{Sensitivity to the
load-bearing threshold. The labels are re-derived from each plan's own wall
thicknesses at each value and the accuracy re-measured. ``Label cols.'' is
the mean number of columns in the resulting label.}
\label{tab:thresholdacc}
\scriptsize
\begin{tabular}{rlrrrrr}
\toprule
Threshold (mm) & Variant & Label cols. & Mean rank & Top-1 & Top-5 & Distance (m) \\
\midrule
\multicolumn{7}{l}{\emph{Fixed weights}} \\
\quad 140 & V1 & 74 & 2854 & 0.000 & 0.000 & 1.060 \\
\quad 140 & V3 & 74 & 2854 & 0.000 & 0.000 & 1.060 \\
\quad 160 & V1 & 57 & 2429 & 0.000 & 0.000 & 1.000 \\
\quad 160 & V3 & 57 & 2429 & 0.000 & 0.000 & 1.000 \\
\quad 180 & V1 & 48 & 2198 & 0.000 & 0.000 & 1.123 \\
\quad 180 & V3 & 48 & 2198 & 0.000 & 0.000 & 1.123 \\
\quad 200 & V1 & 44 & 2162 & 0.000 & 0.000 & 1.159 \\
\quad 200 & V3 & 44 & 2162 & 0.000 & 0.000 & 1.159 \\
\midrule
\multicolumn{7}{l}{\emph{Learned weights}} \\
\quad 140 & V1 & 74 & 24 & 0.143 & 0.429 & 0.208 \\
\quad 140 & V3 & 74 & 24 & 0.143 & 0.429 & 0.232 \\
\quad 160 & V1 & 57 & 22 & 0.000 & 0.429 & 0.248 \\
\quad 160 & V3 & 57 & 22 & 0.143 & 0.429 & 0.271 \\
\quad 180 & V1 & 48 & 90 & 0.000 & 0.143 & 0.402 \\
\quad 180 & V3 & 48 & 90 & 0.000 & 0.143 & 0.431 \\
\quad 200 & V1 & 44 & 194 & 0.000 & 0.143 & 0.568 \\
\quad 200 & V3 & 44 & 195 & 0.000 & 0.143 & 0.612 \\
\bottomrule
\end{tabular}
\end{table}

Under learned weights the result is stable between 140 and 160\,mm, with a
mean rank of about 24 and 22 respectively and distance of 0.232 and
0.271\,m, and degrades beyond, reaching rank 90 at 180\,mm and 195 at
200\,mm with distance rising to 0.612\,m. Part of that degradation is
expected and not independent: the corpus from which the learned weights
were fitted was itself derived at 160\,mm, so the weights are matched to
that threshold. The honest summary is that the headline result tolerates a
$\pm 20$\,mm error in the threshold and is strongly conditioned beyond
that, and this table is the price of the assumption.

\subsection{Score and judge independence}

If the score's ordering and the judge's were nearly identical, the judge
would be adding nothing and variant V1 would be the baseline under another
name. The rank correlation between them, measured across the full candidate
set of each plan, has a mean of 0.161 and a range of $0.076$ to $0.277$.
The judge therefore carries genuinely independent information. That it
nonetheless fails to improve the ranking materially is a separate finding,
and one that the four-variant design exists to expose.
